\documentclass[../main.tex]{subfiles}

\begin{document}
% -----------------------------
\section{Add more operations}
% -----------------------------
As a next step, we will add \texttt{mul} and \texttt{sub} operations and corresponding lowering rules. To do that, first extend dialect:
\begin{lstlisting} [language=Python, caption={Adding new operations in hc\_dialect.py}]
@irdl_op_definition
class HCMul(IRDLOperation):
    name = "hc.mul"
    lhs = operand_def(IntegerType)
    rhs = operand_def(IntegerType)
    res = result_def(IntegerType)

@irdl_op_definition
class HCSub(IRDLOperation):
    name = "hc.sub"
    lhs = operand_def(IntegerType)
    rhs = operand_def(IntegerType)
    res = result_def(IntegerType)
\end{lstlisting}
Also, register new operations:
\begin{lstlisting} [language=Python, caption={Registering new operations in hc\_dialect.py}]
HiCompiler = Dialect(
    "hc",
    (HCAdd, HCMul, HCSub),      # register operations here
    (),  # attrs
)
\end{lstlisting}
In the main file, import operations in order to use them:
\begin{lstlisting} [language=Python, caption={Importing operations}]
from hc_dialect import HiCompiler, HCAdd, HCSub, HCMul
\end{lstlisting}
Modify \path{build_module} function, so that it contains \texttt{HCMul} and \texttt{HCSub}.
\begin{lstlisting}[language=Python, caption={build context snippet}]
c1 = arith.ConstantOp.from_int_and_width(1, 32)
c2 = arith.ConstantOp.from_int_and_width(2, 32)
c3 = arith.ConstantOp.from_int_and_width(3, 32)
c4 = arith.ConstantOp.from_int_and_width(4, 32)

hc_add = HCAdd(
    operands=[c1.result, c2.result],
    result_types=[i32],
)
hc_mul = HCMul(
    operands=[hc_add.results[0], c3.result],
    result_types=[i32],
)
hc_sub = HCSub(
    operands=[hc_mul.results[0], c4.result],
    result_types=[i32],
)

# Return the result
ret = func.ReturnOp(hc_sub.results[0])

# Block containing the constants + hc ops + return
block = Block(ops=[c1, c2, c3, c4, hc_add, hc_mul, hc_sub, ret])
...
\end{lstlisting}
Lastly, modify \path{match_and_rewrite} function:
\begin{lstlisting}[language=Python, caption={build\_context snippet}]
def match_and_rewrite(self, op: Operation, rewriter: PatternRewriter):
    # Debug: uncomment if you want to see traversal
    # print("VISIT:", op.name)

    if op.name not in ("hc.add", "hc.mul", "hc.sub"):
        return

    # operations have two operands
    lhs, rhs = op.operands
    if op.name == "hc.add":
        new_op = arith.AddiOp(lhs, rhs)
    elif op.name == "hc.mul":
        new_op = arith.MuliOp(lhs, rhs)
    elif op.name == "hc.sub":
        new_op = arith.SubiOp(lhs, rhs)
    else:
        raise Exception("Operation not supported")

    rewriter.replace_op(
        op,
        new_ops=[new_op],
        new_results=[new_op.result],
        safe_erase=True,
    )
\end{lstlisting}

\path{apply_lowering} and \texttt{main} functions remain the same. Run the script and you should see
\texttt{hc.add}, \texttt{hc.mul} and \texttt{hc.sub} lowered to
\texttt{arith.addi}, \texttt{arith.muli} and \texttt{arith.subi} respectively.

Our \path{hc_dialect} is a front-end/high level IR. In the following example:

\begin{lstlisting} [language=bash, caption={HC instruction example}]
%2 = hc.add %0, %1 : i32
\end{lstlisting}
\begin{itemize}
    \item We define how to add two values in \texttt{our} language.
    \item Semantics are defined by us, not by the compiler infrastructure.
    \item Compiler doesn't know how to optimize it or emit code for it.
\end{itemize}
After lowering \texttt{hc.add}  →  \texttt{arith.addi} we have:
\begin{itemize}
    \item Precise integer semantics
    \item Defined overflow behavior
    \item Standardized typing rules
    \item Access to existing optimizations
\end{itemize}

\begin{lstlisting} [language=, caption={Typical pipeline}]
hc.*        (our DSL)
 ↓
arith.*     (standard math)
 ↓
scf.*       (loops, control flow)
 ↓
llvm.*      (machine model)
 ↓
x86 / ARM   (real code)
\end{lstlisting}

\end{document}
